\subsection{The exact individual-simple boundary in the iterated three-strand family}

The strand rank is still \(r=3\); \(k\) counts ordinary
tensor factors. Put \(K=\mathbb Q(v)\),
\(\delta=v+v^{-1}\) and \(D=\delta^2\), renaming the
source parameter \(u\) as \(v\).
Theorem~3.4 and Corollary~3.5 of \cite{BlasiakChebyshev},
pp.~12--13, classify the generic simples of
\(K\check H_3^{(k)}\). Besides the two scalar simples,
they are
\begin{equation}\label{eq:simpleEffectivitySourceLabels}
 M_{k,j}=\epsilon_+^{\otimes(k-j)}\otimes\check M^{(j)}
       \simeq\check X^{(k)}(c_{k,j}),
 \qquad
 c_{k,j}=\delta^{2(k-j)}A_j^2,
 \quad A_j=\delta^jT_j(1/\delta),
 \quad 1\le j\le k.
\end{equation}
This list is the published classification, not a new result.
In the source's left-column convention, equation~(13),
p.~8, gives the named generators as
\begin{equation}\label{eq:simpleEffectivitySourceMatrices}
 P_1=\begin{pmatrix}\delta^k&c_{k,j}\\0&0\end{pmatrix},
 \qquad
 P_2=\begin{pmatrix}0&0\\1&\delta^k\end{pmatrix}.
\end{equation}
On the scalar simples their respective common values are
\(\delta^k\) and \(0\), as stated in \S2.4, p.~7.
There is no additional normalization of \(c_{k,j}\).

\begin{finding}[\fid{F63}, exact strict effectivity of the individual simples]
\label{f:iteratedSimpleEffectivityClassification}
For each \(k\ge1\), both scalar simples admit a basis in
which the named generators have entries in
\(\mathbb N[v,v^{-1}]\). The two-dimensional simple
\(M_{k,j}\) admits such a generic basis if and only if
\[
                 j\in\{1,2,3,5\}\cap[1,k].
\]
For every other \(j\), no simultaneous basis even gives
both named generators nonnegative-real Laurent entries.
This is a classification of the named action on individual
generic simples, not of arbitrary direct sums or
source-designated canonical bases.
\end{finding}

\begin{proof}
The elementary Chebyshev recurrence gives
\begin{equation}\label{eq:simpleEffectivityRecurrence}
 A_0=A_1=1,\qquad A_j=2A_{j-1}-DA_{j-2}.
\end{equation}
Thus \(A_j\in\mathbb Z[D]\), and every \(A_j\) is
invariant under \(v\mapsto v^{-1}\).
The following choices are nonzero elements of
\(\mathbb N[v,v^{-1}]\):
\[
\begin{array}{c|c|c}
 j&A_j&B_j=\pm A_j\\ \hline
 1&1&1\\
 2&2-D&v^2+v^{-2}\\
 3&4-3D&3v^2+2+3v^{-2}\\
 5&5D^2-20D+16&5v^4+6+5v^{-4}.
\end{array}
\]
For these indices set
\[
             m=\delta^k,\qquad b=\delta^{k-j}B_j.
\]
Both are nonnegative integral Laurent polynomials and
\(b^2=c_{k,j}\). With \(S=\operatorname{diag}(b,1)\),
the explicit conjugation \(S^{-1}P_iS\) gives
\begin{equation}\label{eq:simpleEffectivityPositiveMatrices}
             P_1=\begin{pmatrix}m&b\\0&0\end{pmatrix},
             \qquad
             P_2=\begin{pmatrix}0&0\\b&m\end{pmatrix}.
\end{equation}
The scalar actions were already effective. This proves
existence for all the asserted cases.

Every other positive index is even with \(j\ge4\), or
odd with \(j\ge7\). In these two cases respectively choose
the odd integer
\[
 \ell=j-1\quad\hbox{or}\quad\ell=j-2.
\]
The indicated bounds on \(j\) give
\[
       \frac{\pi}{3}<\theta=\frac{\ell\pi}{2j}
                       <\frac{\pi}{2}.
\]
Consequently \(x_0=\cos\theta\) satisfies
\(0<x_0<1/2\), and
\(T_j(x_0)=\cos(\ell\pi/2)=0\).
Set
\begin{equation}\label{eq:simpleEffectivityPositiveZero}
 \delta_0=1/x_0>2,\qquad
 v_0=\frac{\delta_0+\sqrt{\delta_0^2-4}}2>1.
\end{equation}
Then \(v_0+v_0^{-1}=\delta_0\), so
\(A_j(v_0)=0\) and \(c_{k,j}(v_0)=0\) for every
\(k\ge j\). On the other hand,
\begin{equation}\label{eq:simpleEffectivityNonzeroAtOne}
       c_{k,j}(1)=2^{2k}\cos^2(j\pi/3)>0.
\end{equation}
Indeed no integer \(j\) makes \(\cos(j\pi/3)\) zero.
Thus \(c_{k,j}\) is a nonzero Laurent polynomial.

Direct multiplication of
\eqref{eq:simpleEffectivitySourceMatrices} gives the
basis-independent identity
\begin{equation}\label{eq:simpleEffectivityTwoWordTrace}
                    \operatorname{Tr}(P_1P_2)=c_{k,j}.
\end{equation}
A nonzero Laurent polynomial with nonnegative real
coefficients is strictly positive at every positive
real parameter. If the two named matrices had such
entries in one basis, their product trace would also
have this property. Equations
\eqref{eq:simpleEffectivityPositiveZero}--\eqref{eq:simpleEffectivityTwoWordTrace}
contradict it. This excludes every remaining index.
\end{proof}

\paragraph{The exact arithmetic consequence.}
For each excluded index,
\(c_{k,j}\in\mathbb Z[v^2,v^{-2}]\), since the
prefactor in \eqref{eq:simpleEffectivitySourceLabels}
is an even power of \(\delta\), and \(A_j\) is a
polynomial in \(D\). Its coefficients have both signs:
its nonzero value and positive real zero exclude all
coefficients being nonnegative, while
\eqref{eq:simpleEffectivityNonzeroAtOne} excludes all
coefficients being nonpositive.

Under either \(v=Q^{n/2}\) or \(v=Q^n\), with \(Q>1\)
a prime power, all distinct visible roots have even
point weights. In a point complex pure of integral
weight \(w\), a root of point weight \(e\) can occur
only in degree \(e-w\), by BBD~5.1.8, pp.~126--127
\cite{BBD}. All its visible Euler multiplicities
would therefore have sign \((-1)^w\). The mixed signs
above and eventual finite power-sum uniqueness give a
contradiction. A fixed root cannot cancel in another
pure degree, and Jordan blocks do not change trace
multiplicities. Thus no bounded uninserted pure point
complex has the specified trace at all sufficiently
large extension degrees. A common Tate twist or
cohomological parity change does not repair it.
Arbitrary mixed complexes and inserted traces remain
outside this exclusion.

\paragraph{What is constructed in the effective cases.}
Let \(\mathcal M,\mathcal B\) be finite graded coefficient
spaces with respective characters \(m,b\) from
\eqref{eq:simpleEffectivityPositiveMatrices}. The two-state
matrix kernels
\[
 \mathcal K_1=
 \begin{pmatrix}\mathcal M&\mathcal B\\0&0\end{pmatrix},
 \qquad
 \mathcal K_2=
 \begin{pmatrix}0&0\\\mathcal B&\mathcal M\end{pmatrix}
\]
have these actual named matrices on split \(K_0\).
Their squares have entries obtained by coefficient
tensor product, and associativity and distributivity
give the explicit quadratic isomorphisms
\[
              \mathcal K_i^2\simeq\mathcal M\otimes\mathcal K_i.
\]
Applying the full-Tate functor of \fid{F61},
\(e\mapsto E(-e)[-2e]\), gives pure weight-zero
derived coefficient kernels under \(v=Q^n\), where
\(\ell\nmid Q\) and \(E/\mathbb Q_\ell\) is finite.
This constructs the named action and quadratic
isomorphisms. It does not construct categorical maps
for every source relation or compatible iteration.

\paragraph{Why the simple and whole-module questions differ.}
The theorem does not classify arbitrary semisimple
direct sums. The actual positive four-dimensional
module in \fid{F66} contains the excluded \(j=4\)
simple; its projector need not be effective.
The entire tensor-output obstructions in \fid{F67}
are separate tests. Nor does the generic balanced
basis identify the actual ambient tensor lattice
at \(v=1\): the \(j=3\) model has the precise
non-\(S_3\) specialization discussed in \fid{F67}.

The positive-zero proof remains valid after the
specific padding factor \(\delta^{2(k-j)}\), since
that factor is strictly positive at \(v_0\).
It makes no claim that a positive polynomial factor
preserves negative coefficients of an arbitrary
whole-tensor trace. No canonical basis labels,
all-strand-rank realization, or refutation of the
sources' signed or factored positivity follows.

% Separate paragraph for Part III.
\paragraph{The \fid{F63} individual-simple classification in Grothendieck's language.}
The actual source modules and explicit conjugations
are affine module/action data over \(\mathbb Q(v)\).
EGA~I, Theorem~1.3.7, p.~86, Corollary~1.3.8, p.~87,
and Theorem~1.4.1, pp.~90--92 provide that dictionary
\cite{EGA}. SGA~3, I, 4.7.2--3 and cote~90, scan~17
provide action language \cite{SGA3,Archive}; they do
not turn an auxiliary positive basis into the
designated canonical basis.

The effective cone \(\mathbb N[v,v^{-1}]\) consists
of actual finite graded representation characters.
SGA~6, V, \S\S2--3 and cote~133, scan~18 supply
the effective and virtual coefficient vocabulary
\cite{SGA6,Archive}, with no \(\lambda\)-ring
assertion for the noncommutative action algebra.
The positive-real-zero argument tests this exact
effective cone. The optional derived point kernels
fit SGA~6, I, Definition~4.7, pp.~121--122 and
Example~4.8, p.~122 \cite{SGA6}; BBD~5.1.5/5.1.8,
pp.~125--127 fixes their integral-weight convention
\cite{BBD}. The specified pure-word arithmetic
comparison remains an additional axiom, absent
from the published simple-module classification.
